\subsection{Three Dimensional Coordinate Systems}

\content{
}{% presentation
\usetikzlibrary{3d}
\hfill\begin{tikzpicture}[x={(-135:.9cm)}, y={(0:1cm)}, z={(90:1cm)}]
\draw[->] (0,0,0) -- (3.4,0,0);
\draw[->] (0,0,0) -- (0,3.4,0);
\draw[->] (0,0,0) -- (0,0,3.4);
\foreach \x/\xlabel in {.5/1,1/2,1.5/3,2/4,2.5/5,3/6} {
    \draw(\x,-.1,0)--(\x,.1,0);
    \node at (\x,0,0) [left=3pt] {\xlabel};
}
\foreach \y/\ylabel in {.5/1,1/2,1.5/3,2/4,2.5/5,3/6} {
    \draw(-.1,\y,0)--(.1,\y,0);
    \node at (0,\y,0) [above=3pt] {\ylabel};
}
\foreach \z/\zlabel in {.5/1,1/2,1.5/3,2/4,2.5/5,3/6} {
    \draw(0,-.1,\z)--(0,.1,\z);
    \node at (0,0,\z) [left=3pt] {\zlabel};
}
\vspace{.5in}
\end{tikzpicture}
}{% nopresentation
}{% comments
We describe 3-dimensional space by picking an origin, and three perpendicular
axis, namely $x$, $y$, and $z$-axis.  these must follow the \textbf{right hand
rule}. 

For any point in 3-dimentional space, these three axis give us a way to describe
where that point is, namely a magnitude in each of the three directions of the
axes. We call this the coordinate of the point. 

\textbf{Give an example of plotting a point in 3-space.}
}

\content{
\subsubsection{Surfaces and Solids}
\begin{example} Describe the surfaces defined by the following equations. 
\begin{itemize}
    \item $z=3$
    \item $y=5$
\end{itemize}
\end{example}
}{% presentation
\usetikzlibrary{3d}
\begin{tikzpicture}[x={(-135:.9cm)}, y={(0:1cm)}, z={(90:1cm)}]
\draw[->] (0,0,0) -- (3.4,0,0);
\draw[->] (0,0,0) -- (0,3.4,0);
\draw[->] (0,0,0) -- (0,0,3.4);
\foreach \x/\xlabel in {.5/1,1/2,1.5/3,2/4,2.5/5,3/6} {
    \draw(\x,-.1,0)--(\x,.1,0);
    \node at (\x,0,0) [left=3pt] {\xlabel};
}
\foreach \y/\ylabel in {.5/1,1/2,1.5/3,2/4,2.5/5,3/6} {
    \draw(-.1,\y,0)--(.1,\y,0);
    \node at (0,\y,0) [above=3pt] {\ylabel};
}
\foreach \z/\zlabel in {.5/1,1/2,1.5/3,2/4,2.5/5,3/6} {
    \draw(0,-.1,\z)--(0,.1,\z);
    \node at (0,0,\z) [left=3pt] {\zlabel};
}
\vspace{.5in}
\end{tikzpicture}
% second diagram for the other problem
\usetikzlibrary{3d}
\hspace{.5in}\begin{tikzpicture}[x={(-135:.9cm)}, y={(0:1cm)}, z={(90:1cm)}]
\draw[->] (0,0,0) -- (3.4,0,0);
\draw[->] (0,0,0) -- (0,3.4,0);
\draw[->] (0,0,0) -- (0,0,3.4);
\foreach \x/\xlabel in {.5/1,1/2,1.5/3,2/4,2.5/5,3/6} {
    \draw(\x,-.1,0)--(\x,.1,0);
    \node at (\x,0,0) [left=3pt] {\xlabel};
}
\foreach \y/\ylabel in {.5/1,1/2,1.5/3,2/4,2.5/5,3/6} {
    \draw(-.1,\y,0)--(.1,\y,0);
    \node at (0,\y,0) [above=3pt] {\ylabel};
}
\foreach \z/\zlabel in {.5/1,1/2,1.5/3,2/4,2.5/5,3/6} {
    \draw(0,-.1,\z)--(0,.1,\z);
    \node at (0,0,\z) [left=3pt] {\zlabel};
}
\vspace{.5in}
\end{tikzpicture}
}{% nopresentation
}{% comments
Mention that in 2-dimensional geometry, graphs of equations involving $x$ and
$y$ are curves in the plane, while here, graphs of equations involving $x$, $y$,
and $z$ are surfaces. 
}

\content{
\begin{example} Describe the set of points which satisfy the two equations
$$ x^2+y^2=1 \quad \text{and}\quad z=3. $$
\end{example}
}{% presentation
\usetikzlibrary{3d}
\hfill\begin{tikzpicture}[x={(-135:.9cm)}, y={(0:1cm)}, z={(90:1cm)}]
\draw[->] (0,0,0) -- (3.4,0,0);
\draw[->] (0,0,0) -- (0,3.4,0);
\draw[->] (0,0,0) -- (0,0,3.4);
\foreach \x/\xlabel in {.5/1,1/2,1.5/3,2/4,2.5/5,3/6} {
    \draw(\x,-.1,0)--(\x,.1,0);
    \node at (\x,0,0) [left=3pt] {\xlabel};
}
\foreach \y/\ylabel in {.5/1,1/2,1.5/3,2/4,2.5/5,3/6} {
    \draw(-.1,\y,0)--(.1,\y,0);
    \node at (0,\y,0) [above=3pt] {\ylabel};
}
\foreach \z/\zlabel in {.5/1,1/2,1.5/3,2/4,2.5/5,3/6} {
    \draw(0,-.1,\z)--(0,.1,\z);
    \node at (0,0,\z) [left=3pt] {\zlabel};
}
\vspace{.5in}
\end{tikzpicture}
}{% nopresentation
}{% comments
If the restriction that $z=3$ had not been given, we would have a cylinder.

Draw out what this would look like with no restriction, and with the restriction
$1\leq z\leq 4$.
}

\content{
\begin{definition} (Distance Formula) Given two points $P_1(x_1,y_1,z_1)$ and 
$P_2(x_2,y_2,z_2)$, the distance between them is 
$$ d(P_1,P_2) = \sqrt{(x_1-x_2)^2+ (y_1-y_2)^2+(z_1-z_2)^2}. $$
\end{definition}

\begin{example} Find the distance between the two points $P_1(2,-1,7)$ and
$P_2(1,-3,5)$.
\end{example}
}{% presentation
\vspace{1in}
}{% nopresentation
}{% comments
}

\content{
\begin{definition} (Equation of a Sphere) An equation of the sphere with center
$P(a,b,c)$ and radius $r\geq 0$ is given by 
$$ (x-a)^2 + (y-b)^2 + (z-c)^2  = r^2 . $$
\end{definition}
}{% presentation
}{% nopresentation
}{% comments
}

\content{
\begin{example} Find the equation of the sphere with center $(1,1,1)$ which passes
through $(2,-1,2)$.
\end{example}
}{% presentation
\vspace{1in}
}{% nopresentation
}{% comments
}

\content{
\begin{example} Show that the equation 
$$ x^2+y^2+z^2 + 4x-6y = 0 $$
is the equation of a sphere, and find it's center and radius. 
\end{example}
}{% presentation
\vspace{2in}
\newpage % end the page here, possibly for a video split. 
}{% nopresentation
}{% comments
}
